\documentclass[a4paper,11pt]{article}
\usepackage[utf8]{inputenc}
\usepackage[T1]{fontenc}
\usepackage[left=0.78in,right=0.78in,top=0.65in,bottom=0.65in]{geometry}
\usepackage{lmodern}
\usepackage{microtype}
\usepackage[hidelinks]{hyperref}
\input{glyphtounicode}
\pdfgentounicode=1
\hypersetup{pdftitle={Peng Li -- Application -- CorPower Ocean -- Hydrodynamics Analysis Engineer},pdfauthor={Peng Li}}
\pagestyle{empty}
\setlength{\parindent}{0pt}
\setlength{\parskip}{8pt}
\setlength{\emergencystretch}{2em}
\urlstyle{same}
\begin{document}
{\LARGE\textbf{Peng Li}}\\[3pt]
Stockholm, Sweden \quad | \quad +46 79 010 31 96\\
\href{mailto:pengli.2810@gmail.com}{pengli.2810@gmail.com}
\hfill 6 November 2026

William Smati and the Recruitment Team\\
CorPower Ocean\\
Stockholm, Sweden

\textbf{Application: Hydrodynamics Analysis Engineer -- Uncrewed Surface Vessel}

Dear William Smati and Recruitment Team,

I am applying for the Hydrodynamics Analysis Engineer position at CorPower Ocean.
I bring experience in computational fluid dynamics (CFD), fluid--structure interaction (FSI),
and structural dynamics, together with industrial model development and validation.
I would like to apply these strengths to understanding vessel response and supporting design
decisions for your uncrewed installation vessel. I am currently completing my Ph.D. research
at KTH Royal Institute of Technology, with my defence scheduled for 4 December 2026.

My doctoral work in upper-airway aeroacoustics predicts how flexible structures respond to
unsteady fluid loading. I developed three-dimensional STAR-CCM+ models coupling compressible LES
with FEM-based structural mechanics, using strong two-way FSI, hyperelastic materials,
and moving overset meshes. By comparing flow conditions, material stiffness, and geometry
variants, I investigate how loads and structural properties affect deformation and vibration.
Modal analysis, cross-spectra, and force--velocity phase relationships help me explain the
mechanisms behind the predicted response. These are transferable skills for analysing
coupled fluid--mechanical systems and comparing engineering design alternatives.

At Joyson Safety Systems, I developed Python-based transient models of pyrotechnical
initiators and coupled their outputs to OpenFOAM simulations of shock-wave propagation
in airbag inflators. I compared predicted pressure histories with measurements and used
parameter-sensitivity studies to assess the influence of component characteristics on
performance and structural integrity. This experience taught me to connect simulation
results with physical behaviour and practical product-development questions.

Verification is central to my modelling approach. I conduct mesh and time-step convergence
studies, compare analytical predictions with finite-element calculations, and use NAISS HPC
resources at PDC and NSC for demanding simulations. My additional research includes
analytical plate bending and buckling, COMSOL modelling of low-Mach-number flows,
and validation of reduced combustion models. A master's degree in Mechanics from
TU Darmstadt, with a final grade of 1.33, provides a foundation in fluid and structural
mechanics and numerical methods for partial differential equations.

CorPower's combination of modelling, vessel development, and feedback from offshore
operations appeals to me because it connects analysis directly to engineering decisions.
My current simulation tools are STAR-CCM+, OpenFOAM, COMSOL, and ABAQUS;
I would build proficiency in OrcaFlex and marine installation analysis as I extend my
CFD--FSI background into this field. I would contribute scientific programming, systematic
model verification, and experience collaborating across fluid mechanics, structural dynamics,
and acoustics. I use English for scientific writing and research communication, and have
led analysis and manuscript preparation for peer-reviewed publications.

I would welcome the opportunity to discuss how my experience in coupled dynamics,
industrial simulation, and numerical model development could contribute to your Marine
Operations team and the deployment of wave-energy technology.

Kind regards,\\[3pt]
\textbf{Peng Li}
\end{document}
